\documentclass{book}
\usepackage{blindtext}

\usepackage[hyphens]{url} % 允许网址在连字符"-"的位置自动换行
\urlstyle{same} % 网址使用与正文相同的字体

\usepackage{hyperref} % 创建超链接
\hypersetup{
    colorlinks=true, % 链接被着色, 默认是false, 黑色字体带有色边框
    linkcolor=blue, % 内部(目录)链接显示为蓝色
    filecolor=magenta, % 本地文件链接为洋红色
    urlcolor=cyan % 网址链接为青色
}

%% BibLaTeX核心命令
% authoryear: 作者-年份  numeric: 数字编号(BibTeX 默认风格)  alphabetic: 字母编号
% apa: 美国心理学会  ieee: 电气电子工程师协会  gb7714-2015numeric: 国标 2015 版
\usepackage[
    backend=biber,    % 编译引擎（必须）
    style=authoryear, % 作者-年份样式（关键）
    natbib=true       % 可选：兼容natbib的\citet/\citep命令
]{biblatex}
\addbibresource{reference.bib} % 必须写.bib后缀

\begin{document}
\frontmatter
\tableofcontents
\clearpage
\addcontentsline{toc}{chapter}{Foreword}
{\huge{\bf Foreword}}

This is foreword.
\clearpage

\mainmatter
\chapter{First Chapter}
This is chapter 1.
\clearpage

\section{First Section} \label{second}
This is section 1.1.

\section{Second Section}
% 网址链接
This is the website of open-source GitHub 
repository: \href{https://github.com/MercuryHaipeng/LaTex-Academic-Writing}{LaTex-Academic-Writing} or go to the next url: \url{https://github.com/MercuryHaipeng/LaTex-Academic-Writing}.

% 本地文件链接，注意链接的文件路径为相对当前LaTex 编译输出 PDF 文件的路径
This is the text of open-source local repository: \href{run:../Chinese-Document.pdf}{LaTeX-Chinese}. 


\section{Reference Management}
%% LaTeX中有两种管理参考文献的方法：一是在 .tex 文档中嵌入参考文献；另一种是使用 BibTeX 进行文献管理。
%% BibTeX文件对不同类型文献的分类如下：
% @article: 期刊论文  @book: 书籍  @inbook: 书中的一部分或某一章节
% @inproceedings/@conference: 会议论文  @mastersthesis/@phdthesis: 硕士/博士论文  
% @proceedings: 会议论文集  @techreport: 技术报告  @misc: 其他类型的文献
%% 丰富的引用命令：
% - \autocite{key}（上下文自适应引用）
% - \textcite{key}（正文式引用，如 “张三 et al., 2020”）
% - \parencite{key}（括号引用，如 “(张三，2020)”）
Some examples for showing how to use \texttt{the bibliography} enviroment:
% 传统的 nabib 引用命令
\begin{itemize}
    \item Book reference sample: The \LaTeX\ companion book \cite{latexcompanion}.
    \item Paper reference sample: On the electrodynamics of moving bodies \citet{einstein}.
    \item Open-source reference sample: Knuth: Computers and Typesetting \citep{knuthwebsite}
\end{itemize}

% BibLaTeX 引用命令
\begin{itemize}
    \item Book reference sample: The \LaTeX\ companion book \autocite{latexcompanion}.
    \item Paper reference sample: On the electrodynamics of moving bodies \textcite{einstein}.
    \item Open-source reference sample: Knuth: Computers and Typesetting \parencite{knuthwebsite}
\end{itemize}

% 插入参考文献
\printbibliography

\end{document}